\QHeading{第1题　在（　　）里填上适当的数。}
\par\noindent


\PlainItem{$9$千克$5$克$=$$(${\color{answerred}$9.005$}$)$千克}
\par

\PlainItem{$6020$千克$=$$(${\color{answerred}$6.02$}$)$吨}
\PlainItem{$5$千米$30$米$=$$(${\color{answerred}$5.03$}$)$千米}
\PlainItem{$435$克$=$$(${\color{answerred}$0.435$}$)$千克}
\PlainItem{$4.35$吨$=$$(${\color{answerred}$4$}$)$吨$(${\color{answerred}$350$}$)$千克}
\PlainItem{$12$米$4$厘米$=$$(${\color{answerred}$12.04$}$)$米}
\PlainItem{$0.2$米$=$$(${\color{answerred}$2$}$)$分米}
\PlainItem{$3$吨$500$千克$=$$(${\color{answerred}$3.5$}$)$吨}
\par\noindent

\Answer{
\QHeading{计算过程：1. 单位换算}
\begin{enumerate}
\item 9千克5克换算成千克：
\Formula{\begin{aligned}[t]&\phantom{={}}9+5\div1000\\&=9+0.005\\&=9.005\text{（千克）}\end{aligned}}

\item 5千米30米换算成千米：
\Formula{\begin{aligned}[t]&\phantom{={}}5+30\div1000\\&=5+0.03\\&=5.03\text{（千米）}\end{aligned}}

\item 4.35吨中的0.35吨换算成千克：
\Formula{\begin{aligned}[t]&\phantom{={}}0.35\times1000\\&=350\text{（千克）}\end{aligned}}
所以4.35吨是4吨350千克。
\item 0.2米换算成分米：
\Formula{\begin{aligned}[t]&\phantom{={}}0.2\times10\\&=2\text{（分米）}\end{aligned}}

\item 6020千克换算成吨：
\Formula{\begin{aligned}[t]&\phantom{={}}6020\div1000\\&=6.02\text{（吨）}\end{aligned}}

\item 435克换算成千克：
\Formula{\begin{aligned}[t]&\phantom{={}}435\div1000\\&=0.435\text{（千克）}\end{aligned}}

\item 12米4厘米换算成米：
\Formula{\begin{aligned}[t]&\phantom{={}}12+4\div100\\&=12+0.04\\&=12.04\text{（米）}\end{aligned}}

\item 3吨500千克换算成吨：
\Formula{\begin{aligned}[t]&\phantom{={}}3+500\div1000\\&=3+0.5\\&=3.5\text{（吨）}\end{aligned}}

\end{enumerate}

}

\QHeading{第2题　用简便方法计算下面各题。}

\CalcBlock[方法一：简便计算]{1}{\begin{aligned}[t]&\phantom{={}}\color{black}50.75-(4.75+0.88)\\&\color{answerred}=50.75-4.75-0.88\\&\color{answerred}=46-0.88\\&\color{answerred}=45.12\end{aligned}}
\par
\MethodBlock{方法二：按运算顺序直接计算（对照）}{本法只作结果核对；本题要求简便计算，前面的简便解法仍保留。}{\begin{aligned}[t]&\phantom{={}}50.75-(4.75+0.88)\\
&=50.75-5.63\\
&=45.12\end{aligned}}


\CalcBlock[方法一：简便计算]{2}{\begin{aligned}[t]&\phantom{={}}\color{black}6.78+9.99\\&\color{answerred}=6.78+(10-0.01)\\&\color{answerred}=6.78+10-0.01\\&\color{answerred}=16.78-0.01\\&\color{answerred}=16.77\end{aligned}}
\MethodBlock{方法二：拆成整数部分和小数部分}{}{\begin{aligned}[t]&\phantom{={}}6.78+9.99\\
&=6.78+9+0.99\\
&=15.78+0.99\\
&=16.77\end{aligned}}

\MethodBlock{方法三：按相同计数单位计算}{先看成若干个百分之一，按整数加减，最后还原成小数。}{\begin{aligned}[t]&\phantom{={}}6.78+9.99\\
&=(678+999)\div 100\\
&=1677\div 100\\
&=16.77\end{aligned}}

\CalcBlock[方法一：简便计算]{3}{\begin{aligned}[t]&\phantom{={}}\color{black}19.1+1.26+2.74+0.9\\&\color{answerred}=19.1+0.9+1.26+2.74\\&\color{answerred}=(19.1+0.9)+(1.26+2.74)\\&\color{answerred}=20+4\\&\color{answerred}=24\end{aligned}}
\MethodBlock{方法二：交换、结合后先凑整}{保留每个数原来的加减号，先把能凑整的部分放在一起。}{\begin{aligned}[t]&\phantom{={}}19.1+1.26+2.74+0.9\\
&=(19.1+0.9)+1.26+2.74\\
&=20+1.26+2.74\\
&=21.26+2.74\\
&=24\end{aligned}}

\MethodBlock{方法三：按运算顺序直接计算（对照）}{本法只作结果核对；本题要求简便计算，前面的简便解法仍保留。}{\begin{aligned}[t]&\phantom{={}}19.1+1.26+2.74+0.9\\
&=20.36+2.74+0.9\\
&=23.1+0.9\\
&=24\end{aligned}}

\CalcBlock[方法一：简便计算]{4}{\begin{aligned}[t]&\phantom{={}}\color{black}6.8+0.2-6.8+0.2\\&\color{answerred}=6.8-6.8+0.2+0.2\\&\color{answerred}=(6.8-6.8)+(0.2+0.2)\\&\color{answerred}=0+0.4\\&\color{answerred}=0.4\end{aligned}}
\MethodBlock{方法二：交换、结合后先凑整}{保留每个数原来的加减号，先把能凑整的部分放在一起。}{\begin{aligned}[t]&\phantom{={}}6.8+0.2-6.8+0.2\\
&=(6.8-6.8)+0.2+0.2\\
&=0+0.2+0.2\\
&=0.2+0.2\\
&=0.4\end{aligned}}

\MethodBlock{方法三：先处理括号或乘除，再提取相同因数}{用乘法分配律反过来计算。}{\begin{aligned}[t]&\phantom{={}}6.8+0.2-6.8+0.2\\
&=7-6.8+0.2\\
&=0.2+0.2\\
&=0.2\times (1+1)\\
&=0.2\times 2\\
&=0.4\end{aligned}}

\CalcBlock[方法一：简便计算]{5}{\begin{aligned}[t]&\phantom{={}}\color{black}13.76-2.99\\&\color{answerred}=13.76-(3-0.01)\\&\color{answerred}=13.76-3+0.01\\&\color{answerred}=10.76+0.01\\&\color{answerred}=10.77\end{aligned}}
\MethodBlock{方法二：拆成整数部分和小数部分}{}{\begin{aligned}[t]&\phantom{={}}13.76-2.99\\
&=13.76-2-0.99\\
&=11.76-0.99\\
&=10.77\end{aligned}}

\MethodBlock{方法三：按相同计数单位计算}{先看成若干个百分之一，按整数加减，最后还原成小数。}{\begin{aligned}[t]&\phantom{={}}13.76-2.99\\
&=(1376-299)\div 100\\
&=1077\div 100\\
&=10.77\end{aligned}}

\MethodBlock{方法四：补差法}{从减数先补到一个整数，再补到被减数，两段差相加。}{\begin{aligned}[t]&\phantom{={}}13.76-2.99\\
&=(3-2.99)+(13.76-3)\\
&=0.01+10.76\\
&=10.77\end{aligned}}

\CalcBlock[方法一：简便计算]{6}{\begin{aligned}[t]&\phantom{={}}\color{black}14.78-4.68-5.32\\&\color{answerred}=14.78-(4.68+5.32)\\&\color{answerred}=14.78-10\\&\color{answerred}=4.78\end{aligned}}
\MethodBlock{方法二：按运算顺序直接计算（对照）}{本法只作结果核对；本题要求简便计算，前面的简便解法仍保留。}{\begin{aligned}[t]&\phantom{={}}14.78-4.68-5.32\\
&=10.1-5.32\\
&=4.78\end{aligned}}

\par\noindent

\QHeading{第3题　解决问题。}

\Problem{① 小明有123.4元零花钱，小慧的零花钱比小明多24.9元，小慧的零花钱有多少元？}
\par


\MethodTitle{方法一}
\Answer{$123.4+24.9=148.3\text{（元）}$\par 答：小慧的零花钱有148.3元。}

\MethodTitle{方法二：凑25元再减0.1元}
\Answer{24.9元比25元少0.1元。}
\Answer{\Formula{\begin{aligned}[t]&\phantom{={}}123.4+24.9\\
&=123.4+25-0.1\\
&=148.4-0.1\\
&=148.3\text{（元）}\end{aligned}}}
\Answer{答：小慧有148.3元零花钱。}
\Problem{② 四（1）班有52人，参加书法班的有25人，参加绘画班的有20人，两种班都不参加的有10人，两种班都参加的有多少人？}

\MethodTitle{方法一}
\Answer{{\Formula{\begin{aligned}[t]&\phantom{={}}25+20-(52-10)\\&=25+20-42\\&=45-42\\&=3\text{（人）}\end{aligned}}}\par 答：两种班都参加的有3人。}

\MethodTitle{方法二：先求只属于第二类的人数}
\Answer{总人数减去第一类人数，再减去两类都不属于的人数，得到仅属于第二类的人数；再从第二类总人数中扣除。}
\Answer{\Formula{\begin{aligned}[t]&\phantom{={}}20-(52-25-10)\\
&=20-(27-10)\\
&=20-17\\
&=3\text{（人）}\end{aligned}}}
\Answer{答：两种班都参加的有3人。}

